% -*- coding: utf-8 -*-
% !TEX program = xelatex
\documentclass[plain,noanswer,medmath]{../../template/jnuexam}
\usepackage{../../template/kysx}

\begin{document}


\examtitle{2019}{3}


\loadproblems{1}{2019P1}%载入试卷一
\loadproblems{2}{2019P2}%载入试卷二


\makepart{选择题}{1～8小题，每小题4分，共32分}


\useproblem{1}{1}{1}%【同试卷一第一(1)题】


\begin{problem}
已知方程$x^5 - 5x + k = 0$有三个不同的实根，则$k$的取值范围是\pickout{}
\begin{abcd}
\item $(-\infty , - 4)$．
\item $(4, +\infty)$．
\item $(- 4,0)$．
\item $(- 4,4)$．
\end{abcd}
\end{problem}


\useproblem{2}{1}{4}%【同试卷二第一(4)题】


\begin{problem}
若级数$\sum_{n = 1}^{\infty}nu_n$绝对收敛，$\sum_{n = 1}^{\infty}\frac{v_n}{n}$条件收敛，则\pickout{}
\begin{abcd}
\item $\sum_{n = 1}^{\infty}u_n v_n$条件收敛．
\item $\sum_{n = 1}^{\infty}u_n v_n$绝对收敛．
\item $\sum_{n = 1}^{\infty}(u_n + v_n)$收敛．
\item $\sum_{n = 1}^{\infty}(u_n + v_n)$发散．
\end{abcd}
\end{problem}


\useproblem{2}{1}{7}%【同试卷二第一(7)题】


\useproblem{1}{1}{5}%【同试卷一第一(5)题】


\useproblem{1}{1}{7}%【同试卷一第一(7)题】


\useproblem{1}{1}{8}%【同试卷一第一(8)题】


\makepart{填空题}{9～14小题，每小题4分，共24分}


\begin{problem}
$\lim_{n \to \infty} \left(\frac{1}{1 \times 2} + \frac{1}{2 \times 3} + \cdots
 + \frac{1}{n \times(n + 1)} \right)^n =$ \fillin{}．
\end{problem}


\begin{problem}%【类似试卷二第一(2)题】
曲线$y = x\sin x + 2\cos x$（$- \frac{\pi}{2} < x < \frac{3\pi}{2}$）的拐点坐标是 \fillin{}．
\end{problem}


\begin{problem}
已知函数$f(x) = \int_1^x \sqrt{1 + t^4} \dt$，则$\int_0^1 x^2 f(x)\dx =$ \fillin{}．
\end{problem}


\begin{problem}
以$P_A$, $P_B$分别表示$A$, $B$两个商品的价格．设商品$A$的需求函数
$$Q_A = 500 - P_A^2 - P_A P_B + 2P_B^2,$$
则当$P_A=10$, $P_B=20$时，商品$A$的需求量对自身价格弹性$\eta_{AA}\;(\eta_{AA} > 0)=$ \fillin{}．
\end{problem}


\begin{problem}
已知矩阵$A = \begin{pmatrix}
  1 & 0 & -1 \\
  1 & 1 & -1 \\
  0 & 1 & a^2 - 1
\end{pmatrix}$，$b = \begin{pmatrix}
  0 \\
  1 \\
  a
\end{pmatrix}$．若线性方程组$Ax = b$有无穷多解，则$a =$ \fillin{}．
\end{problem}


\useproblem{1}{2}{14}%【同试卷一第二(14)题】


\makepart{解答题}{15～23题，共94分}


\useproblem[points=10]{2}{3}{15}%【同试卷二第三(15)题】


\begin{problem}[points=10]
已知$f(u,v)$具有二阶连续偏导数，且$g(x,y) = xy - f(x + y,x - y)$，求
$$\frac{\pd^2 g}{\pd x^2} + \frac{\pd^2 g}{\pdx\pd y} + \frac{\pd^2 g}{\pd y^2}.$$
\end{problem}


\useproblem[points=10]{2}{3}{17}%【同试卷二第三(17)题】


\useproblem[points=10]{1}{3}{17}%【同试卷一第三(17)题】


\useproblem[points=10]{1}{3}{18}%【同试卷一第三(18)题】


\useproblem[points=11]{2}{3}{22}%【同试卷二第三(22)题】


\useproblem[points=11]{1}{3}{21}%【同试卷一第三(21)题】


\useproblem[points=11]{1}{3}{22}%【同试卷一第三(22)题】


\useproblem[points=11]{1}{3}{23}%【同试卷一第三(23)题】


\end{document}
